% Options for packages loaded elsewhere
\PassOptionsToPackage{unicode}{hyperref}
\PassOptionsToPackage{hyphens}{url}
\documentclass[
  11pt,
  a4paper,
]{article}
\usepackage{xcolor}
\usepackage[margin=1in]{geometry}
\usepackage{amsmath,amssymb}
\setcounter{secnumdepth}{-\maxdimen} % remove section numbering
\usepackage{iftex}
\ifPDFTeX
  \usepackage[T1]{fontenc}
  \usepackage[utf8]{inputenc}
  \usepackage{textcomp} % provide euro and other symbols
\else % if luatex or xetex
  \usepackage{unicode-math} % this also loads fontspec
  \defaultfontfeatures{Scale=MatchLowercase}
  \defaultfontfeatures[\rmfamily]{Ligatures=TeX,Scale=1}
\fi
\usepackage{lmodern}
\ifPDFTeX\else
  % xetex/luatex font selection
\fi
% Use upquote if available, for straight quotes in verbatim environments
\IfFileExists{upquote.sty}{\usepackage{upquote}}{}
\IfFileExists{microtype.sty}{% use microtype if available
  \usepackage[]{microtype}
  \UseMicrotypeSet[protrusion]{basicmath} % disable protrusion for tt fonts
}{}
\makeatletter
\@ifundefined{KOMAClassName}{% if non-KOMA class
  \IfFileExists{parskip.sty}{%
    \usepackage{parskip}
  }{% else
    \setlength{\parindent}{0pt}
    \setlength{\parskip}{6pt plus 2pt minus 1pt}}
}{% if KOMA class
  \KOMAoptions{parskip=half}}
\makeatother
\usepackage{color}
\usepackage{fancyvrb}
\newcommand{\VerbBar}{|}
\newcommand{\VERB}{\Verb[commandchars=\\\{\}]}
\DefineVerbatimEnvironment{Highlighting}{Verbatim}{commandchars=\\\{\}}
% Add ',fontsize=\small' for more characters per line
\newenvironment{Shaded}{}{}
\newcommand{\AlertTok}[1]{\textcolor[rgb]{1.00,0.00,0.00}{\textbf{#1}}}
\newcommand{\AnnotationTok}[1]{\textcolor[rgb]{0.38,0.63,0.69}{\textbf{\textit{#1}}}}
\newcommand{\AttributeTok}[1]{\textcolor[rgb]{0.49,0.56,0.16}{#1}}
\newcommand{\BaseNTok}[1]{\textcolor[rgb]{0.25,0.63,0.44}{#1}}
\newcommand{\BuiltInTok}[1]{\textcolor[rgb]{0.00,0.50,0.00}{#1}}
\newcommand{\CharTok}[1]{\textcolor[rgb]{0.25,0.44,0.63}{#1}}
\newcommand{\CommentTok}[1]{\textcolor[rgb]{0.38,0.63,0.69}{\textit{#1}}}
\newcommand{\CommentVarTok}[1]{\textcolor[rgb]{0.38,0.63,0.69}{\textbf{\textit{#1}}}}
\newcommand{\ConstantTok}[1]{\textcolor[rgb]{0.53,0.00,0.00}{#1}}
\newcommand{\ControlFlowTok}[1]{\textcolor[rgb]{0.00,0.44,0.13}{\textbf{#1}}}
\newcommand{\DataTypeTok}[1]{\textcolor[rgb]{0.56,0.13,0.00}{#1}}
\newcommand{\DecValTok}[1]{\textcolor[rgb]{0.25,0.63,0.44}{#1}}
\newcommand{\DocumentationTok}[1]{\textcolor[rgb]{0.73,0.13,0.13}{\textit{#1}}}
\newcommand{\ErrorTok}[1]{\textcolor[rgb]{1.00,0.00,0.00}{\textbf{#1}}}
\newcommand{\ExtensionTok}[1]{#1}
\newcommand{\FloatTok}[1]{\textcolor[rgb]{0.25,0.63,0.44}{#1}}
\newcommand{\FunctionTok}[1]{\textcolor[rgb]{0.02,0.16,0.49}{#1}}
\newcommand{\ImportTok}[1]{\textcolor[rgb]{0.00,0.50,0.00}{\textbf{#1}}}
\newcommand{\InformationTok}[1]{\textcolor[rgb]{0.38,0.63,0.69}{\textbf{\textit{#1}}}}
\newcommand{\KeywordTok}[1]{\textcolor[rgb]{0.00,0.44,0.13}{\textbf{#1}}}
\newcommand{\NormalTok}[1]{#1}
\newcommand{\OperatorTok}[1]{\textcolor[rgb]{0.40,0.40,0.40}{#1}}
\newcommand{\OtherTok}[1]{\textcolor[rgb]{0.00,0.44,0.13}{#1}}
\newcommand{\PreprocessorTok}[1]{\textcolor[rgb]{0.74,0.48,0.00}{#1}}
\newcommand{\RegionMarkerTok}[1]{#1}
\newcommand{\SpecialCharTok}[1]{\textcolor[rgb]{0.25,0.44,0.63}{#1}}
\newcommand{\SpecialStringTok}[1]{\textcolor[rgb]{0.73,0.40,0.53}{#1}}
\newcommand{\StringTok}[1]{\textcolor[rgb]{0.25,0.44,0.63}{#1}}
\newcommand{\VariableTok}[1]{\textcolor[rgb]{0.10,0.09,0.49}{#1}}
\newcommand{\VerbatimStringTok}[1]{\textcolor[rgb]{0.25,0.44,0.63}{#1}}
\newcommand{\WarningTok}[1]{\textcolor[rgb]{0.38,0.63,0.69}{\textbf{\textit{#1}}}}
\usepackage{longtable,booktabs,array}
\usepackage{caption}
\captionsetup[table]{skip=6pt}
\newcounter{none} % for unnumbered tables
\usepackage{calc} % for calculating minipage widths
% Correct order of tables after \paragraph or \subparagraph
\usepackage{etoolbox}
\makeatletter
\patchcmd\longtable{\par}{\if@noskipsec\mbox{}\fi\par}{}{}
\makeatother
\setlength{\emergencystretch}{3em} % prevent overfull lines
\providecommand{\tightlist}{%
  \setlength{\itemsep}{0pt}\setlength{\parskip}{0pt}}
\usepackage{bookmark}
\IfFileExists{xurl.sty}{\usepackage{xurl}}{} % add URL line breaks if available
\urlstyle{same}
% fallback for those not using the hyperref driver hyperxmp:
\makeatletter
\@ifundefined{xmpquote}{\newcommand{\xmpquote}[1]{#1}}{}
\makeatother
\hypersetup{
  hidelinks,
  pdfcreator={LaTeX via pandoc}}

\author{}
\date{}

\begin{document}

\section{FPGA JTAG Programmer}\label{fpga-jtag-programmer}

\textbf{Project report}

October 2026

\subsection{Abstract}\label{abstract}

This project turns a general-purpose FPGA board into a JTAG master.
Embedded software parses a Serial Vector Format (SVF) file with
\texttt{libxsvf} and bit-bangs the target Test Access Port through an
AXI GPIO block. There is no JTAG cable IP and no vendor programming pod
in the data path.

Two programmer platforms share the same player design. A Digilent Nexys
A7 runs the player on MicroBlaze and drives Pmod JA. A Xilinx ZCU102
runs the same player on a Cortex-A53 and drives Pmod J55. The checked-in
SVF configures a second Nexys A7-100T (Artix-7 XC7A100T). A companion
verifier on a Nexys A7 captures the ZCU102 waveforms and checks them
against the same SVF.

A completed configuration reports \texttt{rc=0}, a non-zero TDO check
count, and a lit DONE LED on the target. Playback is slower than a
USB-JTAG pod because each JTAG clock is several AXI GPIO writes from C.

\subsection{1. Introduction}\label{introduction}

Configuring an FPGA over JTAG normally requires a vendor cable and host
software. This project replaces that path with a board that is already
on the bench. The programmer FPGA runs a small standalone application.
The application reads an SVF file that was recorded from a bitstream,
and it drives TCK, TMS, and TDI while sampling TDO.

The programmer board and the target board are different boards. The SVF
configures the target. It does not reconfigure the FPGA that is
executing the player. Programming the same device that is running the
player would replace the player and stop the transfer.

The current asset, \texttt{assets/nexys\_a7\_01.svf}, is a Vivado
configuration SVF for one device:

\begin{itemize}
\tightlist
\item
  The instruction length is 6 bits, which matches a 7-series TAP.
\item
  Chain padding is zero (\texttt{HIR\ 0}, \texttt{TIR\ 0},
  \texttt{HDR\ 0}, \texttt{TDR\ 0}), so the Artix-7 must be the only
  device on the scan chain.
\item
  The IDCODE check expects \texttt{0x03631093} with mask
  \texttt{0x0FFFFFFF}. That is the XC7A100T on a Nexys A7-100T. The mask
  ignores the 4-bit revision field.
\item
  The configuration payload is one data shift of 30,606,304 bits,
  followed by the startup and status checks Vivado emits.
\end{itemize}

\subsection{2. Objectives}\label{objectives}

\begin{enumerate}
\def\labelenumi{\arabic{enumi}.}
\tightlist
\item
  Drive a standard 4-wire JTAG port from an FPGA GPIO block under
  software control.
\item
  Play a vendor-generated SVF, including a full Artix-7 configuration,
  without a programming cable in the data path.
\item
  Support the same player on a soft-core MicroBlaze system and on a Zynq
  UltraScale+ Cortex-A53.
\item
  Provide bring-up commands that exercise the pins before a target is
  attached.
\item
  Provide a second board that can capture the programmer's waveforms and
  compare them with the SVF.
\end{enumerate}

\subsection{3. Background}\label{background}

\subsubsection{3.1 JTAG}\label{jtag}

IEEE 1149.1 defines a Test Access Port of four required signals:

{\def\LTcaptype{none} % do not increment counter
\begin{longtable}[]{@{}lll@{}}
\toprule\noalign{}
Signal & Direction from the master & Role \\
\midrule\noalign{}
\endhead
\bottomrule\noalign{}
\endlastfoot
TCK & Output & Clock \\
TMS & Output & TAP state \\
TDI & Output & Data into the target \\
TDO & Input & Data out of the target \\
\end{longtable}
}

An optional TRST pin is not used. The 7-series TAP is reset by five or
more TCK pulses with TMS held high. The SVF \texttt{STATE\ RESET}
sequence already generates that pattern. \texttt{TRST\ OFF} is accepted
by the player and does not move a wire.

The target TAP on a Nexys A7 is header J10. The same TAP is also reached
by the onboard USB-JTAG circuit on J6. Two masters on TCK, TMS, and TDI
will contend, so the target's PROG USB cable stays unplugged for the
whole SVF run.

\subsubsection{3.2 SVF}\label{svf}

Serial Vector Format is a text record of JTAG operations: state moves,
instruction shifts (\texttt{SIR}), data shifts (\texttt{SDR}), run-test
delays, and expected TDO values with masks. \texttt{libxsvf} parses that
text. It does not know about MicroBlaze, Cortex-A53, AXI, or Pmod pins.
The application supplies a host: a byte source for the SVF, a function
that pulses TCK, and a heap allocator.

The player does not read a bitstream. It plays an SVF that already
contains the JTAG operations for one configuration. Changing the
bitstream means generating a new SVF, replacing the asset, and
rebuilding the ELF. There is no runtime download of the file.

\subsection{4. System architecture}\label{system-architecture}

\begin{Shaded}
\begin{Highlighting}[]
\NormalTok{Host PC                         Programmer board                    Target board}
\NormalTok{Vivado  {-}{-} bitstream / SVF {-}{-}\textgreater{}  CPU + DDR + AXI GPIO  {-}{-} TCK  {-}{-}\textgreater{}   Nexys A7 J10}
\NormalTok{Vitis   {-}{-} ELF with SVF {-}{-}{-}{-}{-}\textgreater{}  UART console          {-}{-} TMS  {-}{-}\textgreater{}   XC7A100T TAP}
\NormalTok{                                Pmod pins             {-}{-} TDI  {-}{-}\textgreater{}}
\NormalTok{                                                      \textless{}{-}{-} TDO {-}{-}}
\NormalTok{                                                      {-}{-}{-} GND {-}{-}{-}}
\end{Highlighting}
\end{Shaded}

Vivado records a program operation into an SVF. The SVF is linked into
the player ELF. At run time the CPU parses the file and writes a GPIO
data register. The Pmod pins are the only connection to the target TAP.

{\def\LTcaptype{none} % do not increment counter
\begin{longtable}[]{@{}
  >{\raggedright\arraybackslash}p{(\linewidth - 4\tabcolsep) * \real{0.3333}}
  >{\raggedright\arraybackslash}p{(\linewidth - 4\tabcolsep) * \real{0.3333}}
  >{\raggedright\arraybackslash}p{(\linewidth - 4\tabcolsep) * \real{0.3333}}@{}}
\toprule\noalign{}
\begin{minipage}[b]{\linewidth}\raggedright
\end{minipage} & \begin{minipage}[b]{\linewidth}\raggedright
Nexys A7 programmer
\end{minipage} & \begin{minipage}[b]{\linewidth}\raggedright
ZCU102 programmer
\end{minipage} \\ 
\midrule\noalign{}
\endhead
\bottomrule\noalign{}
\endlastfoot
Project & \texttt{hw/nexys\_a7\_jtag\_programmer} and
\texttt{hw/nexys\_a7\_jtag\_programmer\_2025.01} &
\texttt{hw/zcu102\_jtag\_programmer} \\
CPU & MicroBlaze & Cortex-A53, standalone domain \\
Application & \texttt{sw/nexys\_a7/app\_mb\_svf\_player} &
\texttt{sw/zcu102/app\_svf\_player} \\
JTAG connector & Pmod JA & Pmod J55 \\
Program memory & 128 MiB DDR2 at \texttt{0x80000000} & PS DDR
(\texttt{psu\_ddr\_0}) \\
Console & Board USB-UART (AXI UART Lite), preset 9600 8N1 & PS UART,
115200 8N1 \\
\end{longtable}
}

The player source is the same design on both platforms: dual-channel AXI
GPIO, channel 1 outputs, channel 2 inputs, and the same bit assignment
(bit 0 TCK, bit 1 TMS, bit 2 TDI, input bit 0 TDO). The ZCU102 binary
still prints the Nexys pin names (\texttt{JA1}, \texttt{JA2},
\texttt{JA3}, \texttt{JA7}). On that board those names mean the J55
pins.

A third application, \texttt{sw/nexys\_a7/app\_mb\_jtag\_verifier}, runs
on a Nexys A7 and acts as the JTAG target for the ZCU102 player. It
walks the same SVF with \texttt{libxsvf} and compares the wires against
the parser's expected bits.

\subsection{5. Hardware design}\label{hardware-design}

\subsubsection{5.1 Nexys A7 programmer}\label{nexys-a7-programmer}

The Nexys A7 programmer is a MicroBlaze system with DDR2, a UART, a
board GPIO for LEDs and switches, and a second GPIO used only as the
JTAG port.

The block design uses these IP cores: MicroBlaze 11.0, MIG 7-series 4.2,
Clocking Wizard 6.0, two AXI GPIO 2.0 cores, AXI UART Lite 2.0,
SmartConnect 1.0, MDM 3.2, and local-memory LMB with a 128 KiB BRAM.
MicroBlaze is built with the debug module, a data-side AXI interface,
and 32 KiB instruction and data caches.

The board clock is \texttt{clk\_100MHz}. The clock wizard produces the
system clock for MicroBlaze, the LMB, both GPIO cores, the UART, and one
SmartConnect. A second wizard output, requested at 200 MHz, is the MIG
reference clock. The MIG user clock has its own reset, and a second
SmartConnect sits on that domain so the instruction and data caches can
reach DDR.

MicroBlaze data space:

{\def\LTcaptype{none} % do not increment counter
\begin{longtable}[]{@{}lll@{}}
\toprule\noalign{}
Peripheral & Base & Range \\
\midrule\noalign{}
\endhead
\bottomrule\noalign{}
\endlastfoot
LMB BRAM & \texttt{0x00000000} & 128 KiB \\
\texttt{axi\_gpio\_0} (JTAG) & \texttt{0x40000000} & 64 KiB \\
\texttt{axi\_gpio\_1} (LEDs, switches) & \texttt{0x40010000} & 64 KiB \\
\texttt{axi\_uartlite\_0} & \texttt{0x40600000} & 64 KiB \\
MIG DDR2 & \texttt{0x80000000} & 128 MiB \\
\end{longtable}
}

Instruction space maps the same LMB and the same DDR window. The linker
script places the application, including the embedded SVF, in the DDR
window. Vectors stay in the local BRAM.

\texttt{axi\_gpio\_0} is the programmer port. Both channels are 4 bits
wide. Channel 1 is all outputs and drives board port \texttt{ck\_a0}.
Channel 2 is wired from input port \texttt{ck\_b0}. The firmware also
programs the direction registers: channel 1 all outputs, channel 2 all
inputs. \texttt{axi\_gpio\_1} is the board LED and switch interface and
is not used by the SVF player.

Pmod JA pinout:

{\def\LTcaptype{none} % do not increment counter
\begin{longtable}[]{@{}lllll@{}}
\toprule\noalign{}
JA pin & FPGA pin & GPIO & JTAG & Direction \\
\midrule\noalign{}
\endhead
\bottomrule\noalign{}
\endlastfoot
JA1 & C17 & channel 1, bit 0 & TCK & output \\
JA2 & D18 & channel 1, bit 1 & TMS & output \\
JA3 & E18 & channel 1, bit 2 & TDI & output \\
JA4 & G17 & channel 1, bit 3 & unused & output \\
JA7 & D17 & channel 2, bit 0 & TDO & input \\
JA8 & E17 & channel 2, bit 1 & unused & input \\
JA9 & F18 & channel 2, bit 2 & unused & input \\
JA10 & G18 & channel 2, bit 3 & unused & input \\
\end{longtable}
}

Pins 5 and 11 are ground. Pins 6 and 12 are 3.3 V and are not connected
to the target.

\subsubsection{5.2 ZCU102 programmer}\label{zcu102-programmer}

The ZCU102 constraints place the same 4-bit output and 4-bit input buses
on Pmod J55, all LVCMOS33. The processing system supplies the Cortex-A53
clock, DDR, and the 115200 baud console. The player looks up
\texttt{XPAR\_XGPIO\_0\_BASEADDR} and uses channel 1 as the output
register and channel 2 as the input register.

{\def\LTcaptype{none} % do not increment counter
\begin{longtable}[]{@{}llll@{}}
\toprule\noalign{}
J55 pin & FPGA pin & JTAG & Direction \\
\midrule\noalign{}
\endhead
\bottomrule\noalign{}
\endlastfoot
J55.1 & A20 & TCK & output \\
J55.3 & B20 & TMS & output \\
J55.5 & A22 & TDI & output \\
J55.2 & B21 & unused & output \\
J55.7 & A21 & TDO & input \\
J55.4 & C21 & unused & input \\
J55.6 & C22 & unused & input \\
J55.8 & D21 & unused & input \\
\end{longtable}
}

J55 is on bank 47, powered by \texttt{VADJ\_FPGA}. The constraints use
LVCMOS33, so \texttt{VADJ\_FPGA} is set to 3.3 V before the header is
connected to a Nexys A7.

\subsubsection{5.3 Target}\label{target}

The target is a second Nexys A7. Header J10 is a single-device TAP,
which matches an SVF generated with no header or trailer bits.

\begin{Shaded}
\begin{Highlighting}[]
\NormalTok{Programmer TCK  {-}\textgreater{}  target J10 TCK}
\NormalTok{Programmer TMS  {-}\textgreater{}  target J10 TMS}
\NormalTok{Programmer TDI  {-}\textgreater{}  target J10 TDI}
\NormalTok{Programmer TDO  \textless{}{-}  target J10 TDO}
\NormalTok{Programmer GND \textless{}{-}\textgreater{}  target J10 GND}
\end{Highlighting}
\end{Shaded}

The J10 VREF pin is left open. The programmer does not source or load
that pin. JP3 selects the barrel jack, the PROG USB plug is out, and JP1
is in the JTAG position so QSPI, microSD, or USB cannot load a different
bitstream after the SVF finishes. The target is powered and out of reset
before playback starts. The DONE LED lights after a successful
configuration.

Pmod JA on the Nexys A7 and J55 in this ZCU102 design are LVCMOS33. The
Nexys A7 JTAG bank is 3.3 V, so a Nexys-to-Nexys or ZCU102-to-Nexys
connection does not need level shifters. Any other target voltage does.
The grounds are always connected. The programmer's 3.3 V pin is never
tied to the target.

\subsubsection{5.4 Repository layout}\label{repository-layout}

\begin{Shaded}
\begin{Highlighting}[]
\NormalTok{hw/nexys\_a7\_jtag\_programmer/             Vivado project}
\NormalTok{hw/nexys\_a7\_jtag\_programmer\_2025.01/     Recreate Tcl for the 2025.1 flow}
\NormalTok{hw/zcu102\_jtag\_programmer/               ZCU102 project and J55 constraints}
\NormalTok{sw/nexys\_a7/app\_mb\_svf\_player/           MicroBlaze SVF player}
\NormalTok{sw/nexys\_a7/app\_mb\_jtag\_verifier/        MicroBlaze JTAG verifier}
\NormalTok{sw/nexys\_a7/ws\_20260819/                 Vitis workspace and pf\_nexys\_a7 platform}
\NormalTok{sw/zcu102/app\_svf\_player/                Cortex{-}A53 SVF player}
\NormalTok{sw/zcu102/app\_hello\_world/               Cortex{-}A53 hello world}
\NormalTok{sw/zcu102/ws\_20260819/                   Vitis workspace and pf\_zynq platform}
\NormalTok{docs/                                    Documentation}
\end{Highlighting}
\end{Shaded}

Vivado project files (\texttt{.xpr}), exported hardware (\texttt{.xsa}),
and bitstreams (\texttt{.bit}) are kept in version control. Generated
implementation directories, Vitis build trees, and ELF files are
ignored.

\subsection{6. Firmware}\label{firmware}

\subsubsection{6.1 Host callbacks}\label{host-callbacks}

\texttt{libxsvf\_play} is called in SVF mode. The application implements
\texttt{struct\ libxsvf\_host}:

{\def\LTcaptype{none} % do not increment counter
\begin{longtable}[]{@{}
  >{\raggedright\arraybackslash}p{(\linewidth - 2\tabcolsep) * \real{0.5000}}
  >{\raggedright\arraybackslash}p{(\linewidth - 2\tabcolsep) * \real{0.5000}}@{}}
\toprule\noalign{}
\begin{minipage}[b]{\linewidth}\raggedright
Callback
\end{minipage} & \begin{minipage}[b]{\linewidth}\raggedright
Behavior
\end{minipage} \\
\midrule\noalign{}
\endhead
\bottomrule\noalign{}
\endlastfoot
\texttt{setup} & Clear counters, park TCK high, set GPIO directions \\
\texttt{getbyte} & Next byte from the active buffer, or \texttt{-1} at
the end \\
\texttt{pulse\_tck} & One JTAG clock, described below \\
\texttt{udelay} & Extra TCK pulses for \texttt{RUNTEST\ n\ TCK}, then a
time delay split into 1 s chunks \\
\texttt{set\_frequency} & Print the requested rate and return success.
The GPIO path is not retimed \\
\texttt{set\_trst} & Accept the command. No TRST pin is assigned \\
\texttt{realloc} & C \texttt{realloc} / \texttt{free} on the application
heap \\
\texttt{report\_status}, \texttt{report\_error} & UART text, truncated
to 96 characters \\
\texttt{report\_device} & Print the IDCODE \texttt{libxsvf} shifted
in \\
\texttt{shutdown} & Outputs low \\
\end{longtable}
}

The build compiles \texttt{svf.c}, \texttt{tap.c}, \texttt{play.c},
\texttt{statename.c}, and \texttt{memname.c} from \texttt{libxsvf}.
\texttt{LIBXSVF\_WITHOUT\_XSVF} and \texttt{LIBXSVF\_WITHOUT\_SCAN}
leave XSVF and scan mode out. Optimization is \texttt{-O2}.

\subsubsection{6.2 TCK waveform}\label{tck-waveform}

The GPIO base is \texttt{XPAR\_XGPIO\_0\_BASEADDR}. Channel 1 is the
output data register at offset \texttt{0x0}. Channel 2 is the input data
register at offset \texttt{0x8}.

{\def\LTcaptype{none} % do not increment counter
\begin{longtable}[]{@{}lll@{}}
\toprule\noalign{}
Bit & Mask & Signal \\
\midrule\noalign{}
\endhead
\bottomrule\noalign{}
\endlastfoot
channel 1, bit 0 & \texttt{0x1} & TCK \\
channel 1, bit 1 & \texttt{0x2} & TMS \\
channel 1, bit 2 & \texttt{0x4} & TDI \\
channel 2, bit 0 & \texttt{0x1} & TDO \\
\end{longtable}
}

\texttt{out\_shadow} holds the last output word. Each callback updates
TMS or TDI in that shadow, then writes the whole nibble. Unrelated bits
stay put.

\texttt{mb\_svf\_pulse\_tck} is called once per JTAG clock.
\texttt{libxsvf} passes TMS, an optional TDI bit (\texttt{-1} means
leave TDI alone), an optional expected TDO bit (\texttt{-1} means do not
check), and a return mask. The pulse is negative, which is the edge
order this port of \texttt{libxsvf} is written for:

\begin{Shaded}
\begin{Highlighting}[]
\NormalTok{TCK high  {-}\textgreater{}  write GPIO}
\NormalTok{TCK low   {-}\textgreater{}  write GPIO}
\NormalTok{TCK high  {-}\textgreater{}  write GPIO}
\NormalTok{sample TDO if a check was requested}
\end{Highlighting}
\end{Shaded}

Setup starts with TCK already high, so the first edge the target sees is
the falling edge. Shutdown drives all four output bits low.

In demo mode each of those three writes is followed by
\texttt{usleep(100000)}. Real playback uses a delay of 0 and runs at the
speed of the GPIO writes.

TDO is sampled after the pulse when the parser asked for a value or set
the return mask. If \texttt{ignore\_tdo} is clear and the sampled bit
differs from the expected bit, the callback records the clock count and
both bit values, then returns \texttt{-1}. \texttt{libxsvf} stops and
reports a TDO mismatch. Demo mode sets \texttt{ignore\_tdo}, so a
floating TDO pin does not fail the run.

\texttt{FREQUENCY\ 1.00E+07\ HZ} in the real SVF is a record of what
Vivado asked for. AXI GPIO writes from C are far slower than 10 MHz.
JTAG configuration still completes; it takes longer than a USB-JTAG pod.

\subsubsection{6.3 Where the SVF lives}\label{where-the-svf-lives}

\texttt{svf\_blob.S} is a template:

\begin{Shaded}
\begin{Highlighting}[]
\NormalTok{.}\BuiltInTok{section} \OperatorTok{.}\NormalTok{rodata}\OperatorTok{.}\NormalTok{svf}\OperatorTok{,} \StringTok{"a"}
\NormalTok{.incbin }\StringTok{"@APP\_SVF\_FILE@"}
\end{Highlighting}
\end{Shaded}

CMake writes \texttt{svf\_blob.generated.S} into the build directory
with \texttt{@APP\_SVF\_FILE@} replaced by the absolute path of
\texttt{assets/nexys\_a7\_01.svf}. The checked-in sources do not contain
a machine-specific path. The C code sees the blob through linker symbols
\texttt{nexys\_a7\_01\_svf\_start} and \texttt{nexys\_a7\_01\_svf\_end}.
The size printed at startup is the difference of those addresses.

On the Nexys A7 the linker maps \texttt{.text} and \texttt{.rodata}
(including \texttt{.rodata.svf}) to DDR at \texttt{0x80000000}, length
128 MiB. On the ZCU102 the same sections map to \texttt{psu\_ddr\_0} at
\texttt{0x0}. On both targets the heap is 32 MiB and the stack is 16
KiB.

The 32 MiB heap is required. \texttt{libxsvf} reads a whole command
before executing it, and the configuration command is an \texttt{SDR} of
30,606,304 bits. A default heap cannot hold that command buffer and the
parsed bit arrays. If a rebuilt SVF is substantially larger,
\texttt{\_HEAP\_SIZE} in \texttt{src/lscript.ld} is raised, and the heap
stays inside the DDR region.

\subsubsection{6.4 UART commands}\label{uart-commands}

The UART menu is the whole user interface. The loop never exits.

{\def\LTcaptype{none} % do not increment counter
\begin{longtable}[]{@{}
  >{\raggedright\arraybackslash}p{(\linewidth - 4\tabcolsep) * \real{0.3333}}
  >{\raggedright\arraybackslash}p{(\linewidth - 4\tabcolsep) * \real{0.3333}}
  >{\raggedright\arraybackslash}p{(\linewidth - 4\tabcolsep) * \real{0.3333}}@{}}
\toprule\noalign{}
\begin{minipage}[b]{\linewidth}\raggedright
Key
\end{minipage} & \begin{minipage}[b]{\linewidth}\raggedright
Mode
\end{minipage} & \begin{minipage}[b]{\linewidth}\raggedright
Target required
\end{minipage} \\
\midrule\noalign{}
\endhead
\bottomrule\noalign{}
\endlastfoot
\texttt{t} & Walk a 4-bit pattern on the output pins
(\texttt{1,\ 2,\ 4,\ 8,\ F,\ 0}, 300 ms per step) & No \\
\texttt{d} & Short SVF, slowed to about 100 ms per edge, TDO ignored &
No \\
\texttt{p} & Embedded configuration SVF at GPIO speed, TDO checked &
Yes \\
\texttt{r} & Drive the outputs low & No \\
\end{longtable}
}

\texttt{t} does not call \texttt{libxsvf}. It is the first check that
the output wiring and the bitstream are the JTAG design. \texttt{d}
plays a short string compiled into the ELF: reset, idle, an IDCODE
instruction, a 32-bit data shift, eight run-test clocks, and a short
\texttt{0x55AA} shift. \texttt{p} is the programming path. Every TDO
expectation in the file is enforced. The first one is the IDCODE shift,
so a missing or wrong target fails immediately and the large
configuration shift never starts.

Phases inside the configuration file:

{\def\LTcaptype{none} % do not increment counter
\begin{longtable}[]{@{}
  >{\raggedright\arraybackslash}p{(\linewidth - 2\tabcolsep) * \real{0.5000}}
  >{\raggedright\arraybackslash}p{(\linewidth - 2\tabcolsep) * \real{0.5000}}@{}}
\toprule\noalign{}
\begin{minipage}[b]{\linewidth}\raggedright
SVF section
\end{minipage} & \begin{minipage}[b]{\linewidth}\raggedright
JTAG effect
\end{minipage} \\
\midrule\noalign{}
\endhead
\bottomrule\noalign{}
\endlastfoot
\texttt{STATE\ RESET} / \texttt{STATE\ IDLE} & TMS sequences the TAP
through Test-Logic-Reset into Run-Test/Idle \\
\texttt{SIR\ 6\ TDI\ (09)} + 32-bit \texttt{SDR} & Shift the IDCODE
instruction and read 32 bits. Expect \texttt{0x03631093} under mask
\texttt{0x0FFFFFFF} \\
\texttt{SIR\ 6\ TDI\ (0b)} & JPROGRAM \\
\texttt{SIR\ 6\ TDI\ (14)} and the \texttt{RUNTEST} pair &
Initialization poll, including a 0.1 s wait and 10,000 clocks \\
\texttt{SIR\ 6\ TDI\ (05)} + \texttt{SDR\ 30606304} & Shift the
configuration bitstream into the device \\
\texttt{config/start} & \texttt{JSTART} (\texttt{0x0C}) and clocks that
complete startup \\
\texttt{config/status} & Status read. One bit is checked under mask
\texttt{0x08000000} \\
\end{longtable}
}

Status lines on the UART are clipped to 96 characters. The configuration
\texttt{SDR} is millions of characters, and printing it would dominate
the runtime.

\subsection{7. SVF generation}\label{svf-generation}

The checked-in file is produced by recording a program operation against
an SVF target in Vivado Hardware Manager. No cable and no board are
required. The part is the Nexys A7-100T device,
\texttt{xc7a100tcsg324-1}.

\begin{Shaded}
\begin{Highlighting}[]
\NormalTok{open\_hw\_manager}
\NormalTok{connect\_hw\_server {-}url localhost}\OperatorTok{:}\DecValTok{3121}
\NormalTok{create\_hw\_target nexys\_a7\_svf}
\NormalTok{open\_hw\_target }\KeywordTok{[lindex} \KeywordTok{[}\NormalTok{get\_hw\_targets {-}regexp .}\OperatorTok{*/}\NormalTok{nexys\_a7\_svf}\KeywordTok{]} \DecValTok{0}\KeywordTok{]}
\KeywordTok{set}\NormalTok{ dev }\KeywordTok{[}\NormalTok{create\_hw\_device {-}part xc7a100tcsg324{-}1}\KeywordTok{]}
\NormalTok{set\_property PROGRAM.FILE }\KeywordTok{\{}\NormalTok{C}\OperatorTok{:/}\NormalTok{path}\OperatorTok{/}\NormalTok{to}\OperatorTok{/}\NormalTok{target\_design.bit}\KeywordTok{\}} \DataTypeTok{$dev}
\NormalTok{program\_hw\_devices }\DataTypeTok{$dev}
\NormalTok{write\_hw\_svf }\KeywordTok{\{}\NormalTok{C}\OperatorTok{:/}\NormalTok{path}\OperatorTok{/}\NormalTok{to}\OperatorTok{/}\NormalTok{nexys\_a7\_01.svf}\KeywordTok{\}}
\NormalTok{close\_hw\_target}
\end{Highlighting}
\end{Shaded}

The header of the resulting file is a single-device Artix-7 sequence:
\texttt{TRST\ OFF}, idle end states, \texttt{STATE\ RESET},
\texttt{STATE\ IDLE}, \texttt{FREQUENCY\ 1.00E+07\ HZ}, zero chain
padding, then the IDCODE, JPROGRAM, initialization, and configuration
shifts. Later comments mark \texttt{config/start} and
\texttt{config/status}.

\texttt{SIR\ 6\ TDI\ (09)} is the 7-series IDCODE instruction.
\texttt{SIR\ 6\ TDI\ (0b)} is JPROGRAM. The 30,606,304-bit \texttt{SDR}
is the configuration payload: 3,825,788 bytes of shift data, stored as
ASCII hex, which is why the file is about 7.7 MB.

Before the file is embedded it is checked for the expected IDCODE, zero
padding, a configuration shift of tens of millions of bits, and a
complete startup and status trailer. Vivado writes the init delay as
\texttt{RUNTEST\ 0.100000\ SEC;}. The bundled \texttt{libxsvf} rejects a
decimal point in that position. The file is rewritten as
\texttt{RUNTEST\ 1E-1\ SEC;}, which is the same 0.1 s.

The same asset is installed in all three applications:

\begin{Shaded}
\begin{Highlighting}[]
\NormalTok{sw/nexys\_a7/app\_mb\_svf\_player/assets/nexys\_a7\_01.svf}
\NormalTok{sw/nexys\_a7/app\_mb\_jtag\_verifier/assets/nexys\_a7\_01.svf}
\NormalTok{sw/zcu102/app\_svf\_player/assets/nexys\_a7\_01.svf}
\end{Highlighting}
\end{Shaded}

Replacing the SVF without rebuilding leaves the previous payload in the
ELF. The size printed at startup is the check.

\subsection{8. Verification}\label{verification}

\subsubsection{8.1 Bring-up without a
target}\label{bring-up-without-a-target}

Before J10 is wired, the output pins can drive LEDs through 330 ohm to 1
kohm resistors. TDO is an input and does not get an LED. The LEDs are
removed before a real target is connected.

{\def\LTcaptype{none} % do not increment counter
\begin{longtable}[]{@{}
  >{\raggedright\arraybackslash}p{(\linewidth - 2\tabcolsep) * \real{0.5000}}
  >{\raggedright\arraybackslash}p{(\linewidth - 2\tabcolsep) * \real{0.5000}}@{}}
\toprule\noalign{}
\begin{minipage}[b]{\linewidth}\raggedright
Command
\end{minipage} & \begin{minipage}[b]{\linewidth}\raggedright
Result with no target
\end{minipage} \\
\midrule\noalign{}
\endhead
\bottomrule\noalign{}
\endlastfoot
\texttt{t} & Output nibble walks \texttt{1,\ 2,\ 4,\ 8,\ F,\ 0} with 300
ms between steps \\
\texttt{d} & Parser and pin callbacks run the short demo SVF. TDO
mismatches are ignored. \texttt{rc} is 0 if the parser accepts the
string \\
\texttt{p} & Fails at the IDCODE \texttt{SDR}. UART shows a TDO
mismatch, then \texttt{rc=-1}, and the first failing TCK count plus the
expected and sampled bits \\
\end{longtable}
}

\texttt{p} without a target is a useful negative test. The parser got
far enough to shift 32 data bits and compare them. It is not a sign that
the SVF file or the GPIO path is broken.

\subsubsection{8.2 Programming a target}\label{programming-a-target}

The programmer is built and downloaded over its own USB-JTAG port. The
console banner ends with the blob size, then the menu waits. \texttt{t}
and \texttt{d} are run with J10 open. Outputs are driven low with
\texttt{r}, LEDs are removed, and the four JTAG wires plus ground are
connected. \texttt{p} is sent only after the target is powered.

Success is all of the following:

\begin{itemize}
\tightlist
\item
  \texttt{SVF\ playback\ finished\ with\ rc=0}
\item
  A non-zero TDO check count, and no TDO mismatch line
\item
  The target DONE LED is on
\item
  The design that was baked into the SVF is the design now running on
  the target
\end{itemize}

A mismatch after a long run, with a large TCK count, means IDCODE and
the init poll passed and a later compare failed. The usual causes are a
target that lost power during the shift, another master on TMS or TCK,
or an SVF generated for a different part. A mismatch on the first data
shift is wiring, power, a second master on J6, a swapped TDI/TDO pair,
or a device that is not an XC7A100T.

The final status shift compares a single masked bit
(\texttt{MASK\ (08000000)}). A mismatch there means the device did not
reach the expected status after startup, so \texttt{rc} is \texttt{-1}
even if the long shift itself completed.

\subsubsection{8.3 Waveform verifier}\label{waveform-verifier}

The Nexys A7 verifier is the counterpart of the ZCU102 player. It does
not implement a TAP state machine of its own and it does not write a
second SVF parser. It runs the same \texttt{libxsvf\_play} over the same
SVF bytes, with \texttt{pulse\_tck} inverted:

{\def\LTcaptype{none} % do not increment counter
\begin{longtable}[]{@{}lll@{}}
\toprule\noalign{}
& ZCU102 player & Nexys A7 verifier \\
\midrule\noalign{}
\endhead
\bottomrule\noalign{}
\endlastfoot
TCK & drives the pulse & waits for the pulse \\
TMS & drives the expected value & samples and compares \\
TDI & drives the expected value & samples and compares \\
TDO & samples and compares & drives the expected value \\
\end{longtable}
}

Because \texttt{libxsvf} is deterministic, both instances walk the file
in lockstep. On every clock the verifier receives the TMS and TDI values
the player should be driving, and the TDO value the SVF says the target
should return. Driving TDO lets the player's own TDO checks pass, which
closes the return path.

The verifier detects TCK edges by polling the AXI GPIO input register
from C. It needs real time between the player's falling edge and the
player's TDO sample. The player's \texttt{d} command, at about 100 ms
per edge, leaves that margin and is the bring-up pair. The player's
\texttt{p} command, with zero edge delay, is faster than the poll loop,
so the embedded configuration file is not verified this way. The demo
SVF has no TDO fields, so a \texttt{d} run validates TCK, TMS, and TDI
only.

Wiring is point to point and is not a straight ribbon. On the verifier,
TCK, TMS, and TDI are inputs on JA7, JA8, and JA9, and TDO is an output
on JA1. The verifier is armed before the player starts clocking. The
first clock waits for TCK high and then for the falling edge, so an
idle-low TCK is not mistaken for the first pulse.

\subsection{9. Results}\label{results}

These are the outcomes the player and the checked-in SVF are built to
produce. A bench log of one particular programming session is not stored
in the repository. The pass and fail text is what the firmware prints,
and the IDCODE and shift length are what
\texttt{assets/nexys\_a7\_01.svf} contains.

{\def\LTcaptype{none} % do not increment counter
\begin{longtable}[]{@{}
  >{\raggedright\arraybackslash}p{(\linewidth - 2\tabcolsep) * \real{0.5000}}
  >{\raggedright\arraybackslash}p{(\linewidth - 2\tabcolsep) * \real{0.5000}}@{}}
\toprule\noalign{}
\begin{minipage}[b]{\linewidth}\raggedright
Check
\end{minipage} & \begin{minipage}[b]{\linewidth}\raggedright
Value
\end{minipage} \\
\midrule\noalign{}
\endhead
\bottomrule\noalign{}
\endlastfoot
Device & One TAP, no chain padding \\
Instruction width & 6 \\
IDCODE & \texttt{0x03631093}, mask \texttt{0x0FFFFFFF} \\
Requested TCK & 10 MHz, not enforced by the player \\
Configuration shift & 30,606,304 bits \\
Embedded file size & 7,682,571 bytes (7,712,737 with CRLF line
endings) \\
Recorded Nexys A7 ELF & \texttt{text} 7,747,420, \texttt{data} 1,624,
\texttt{bss} 33,571,276 \\
\end{longtable}
}

The \texttt{text} size tracks the SVF because \texttt{.incbin} places
the file in \texttt{.rodata}. The \texttt{bss} size tracks the 32 MiB
heap reservation. Both are expected for this asset.

A completed configuration reports:

\begin{Shaded}
\begin{Highlighting}[]
\NormalTok{SVF playback finished with rc=0}
\NormalTok{TCK clocks: \textless{}large count\textgreater{}, TDI bits: \textless{}large count\textgreater{}, TDO checks: \textless{}count\textgreater{}}
\end{Highlighting}
\end{Shaded}

The clock count is dominated by the 30,606,304-bit shift plus the
\texttt{RUNTEST} clocks (10,000 during init, 100,000 at startup, and the
smaller sequences). Runtime is the cost of those GPIO transactions. The
10 MHz \texttt{FREQUENCY} command only produces the UART line
\texttt{using\ GPIO\ bit-bang\ speed.}

\subsection{10. Limitations}\label{limitations}

\begin{itemize}
\tightlist
\item
  TRST is not a pin. TAP reset still happens through TMS and TCK.
\item
  The player will not discover extra devices. Padding is fixed in the
  SVF at zero.
\item
  Command \texttt{p} cannot configure the FPGA that is executing the
  player. That download uses the board's own USB-JTAG, or the usual
  PS/PL programming path on the ZCU102.
\item
  The GPIO path does not honor the SVF frequency. Configuration works
  and takes longer than a cable.
\item
  The polling verifier can follow the slowed demo. It cannot follow a
  full-speed configuration shift.
\item
  The loop-test SVF in the verifier, which would exercise TDO as well as
  TCK, TMS, and TDI, has no matching command on the current ZCU102
  player.
\item
  Replacing the SVF without rebuilding leaves the previous payload in
  the ELF.
\end{itemize}

\subsection{11. Conclusion}\label{conclusion}

A MicroBlaze system on a Nexys A7 and a Cortex-A53 system on a ZCU102
both play a Vivado configuration SVF through an AXI GPIO port and
configure a second Nexys A7-100T. The software contract is small:
\texttt{libxsvf} parses the file, and one callback writes TCK, TMS, and
TDI and samples TDO. The same callback, inverted, lets a third board
check the waveform against the file.

The design is sufficient for a single Artix-7 on a 3.3 V TAP. Bring-up
is separated from programming: \texttt{t} and \texttt{d} prove the pins
with no target, and \texttt{p} enforces every TDO compare, starting with
the IDCODE. The cost of the approach is speed. Each bit is several
register writes, and the 30.6 million bit configuration shift dominates
the run.

\subsection{References}\label{references}

\begin{enumerate}
\def\labelenumi{\arabic{enumi}.}
\tightlist
\item
  IEEE Std 1149.1, IEEE Standard for Test Access Port and Boundary-Scan
  Architecture.
\item
  Clifford Wolf, \texttt{libxsvf},
  https://github.com/cliffordwolf/libxsvf.
\item
  Digilent, Nexys A7 FPGA Board Reference Manual.
\item
  AMD, Vivado Design Suite and Vitis embedded tools, 2025.1 / 2025.2
  flows used by this repository.
\item
  AMD, 7 Series FPGAs Configuration User Guide, for the XC7A100T JTAG
  instructions used by the generated SVF (IDCODE, JPROGRAM, JSTART).
\end{enumerate}

\end{document}
